\chapter{Tests, validation et discussion}

\guided{Ce chapitre est d\'eterminant. Les r\'esultats doivent \^etre analys\'es, pas seulement montr\'es. Le lecteur doit comprendre : ce qui a \'et\'e test\'e, comment, avec quels indicateurs, dans quelles conditions, et avec quel niveau de confiance.}

\section{Plan de validation}
\todo{D\'ecrire les types de tests : unitaires, int\'egration, syst\`eme, terrain, validation exp\'erimentale, comparaison \`a l'\'etat de l'art, etc.}

\section{Protocoles et conditions exp\'erimentales}
\todo{D\'ecrire les jeux de donn\'ees, bancs d'essai, capteurs, sc\'enarios, param\`etres et indicateurs.}

\section{Pr\'esentation des r\'esultats}
\begin{table}[H]
\centering
\caption{Exemple de synth\`ese des tests}
\begin{tabularx}{\textwidth}{>{\RaggedRight\arraybackslash}p{2cm} >{\RaggedRight\arraybackslash}X >{\Centering\arraybackslash}p{2.3cm} >{\Centering\arraybackslash}p{2.3cm} >{\RaggedRight\arraybackslash}p{2.2cm}}
\toprule
Test & Objectif & Attendu & Obtenu & Conclusion \\
\midrule
T1 & \todo{...} & \todo{...} & \todo{...} & \todo{Valid\'e / \`a am\'eliorer} \\
T2 & \todo{...} & \todo{...} & \todo{...} & \todo{Valid\'e / \`a am\'eliorer} \\
T3 & \todo{...} & \todo{...} & \todo{...} & \todo{Valid\'e / \`a am\'eliorer} \\
\bottomrule
\end{tabularx}
\end{table}

\section{Discussion critique}
\todo{Interpr\'eter les r\'esultats, expliquer les \'ecarts, discuter les limites, menaces sur la validit\'e et pistes d'am\'elioration.}

\guided{\disciplinehint}
